%%=============================================================================
%% Methodologie
%%=============================================================================

\chapter{\IfLanguageName{dutch}{Methodologie}{Methodology}}%
\label{ch:methodologie}

%% TODO: In dit hoofstuk geef je een korte toelichting over hoe je te werk bent
%% gegaan. Verdeel je onderzoek in grote fasen, en licht in elke fase toe wat
%% de doelstelling was, welke deliverables daar uit gekomen zijn, en welke
%% onderzoeksmethoden je daarbij toegepast hebt. Verantwoord waarom je
%% op deze manier te werk gegaan bent.
%% 
%% Voorbeelden van zulke fasen zijn: literatuurstudie, opstellen van een
%% requirements-analyse, opstellen long-list (bij vergelijkende studie),
%% selectie van geschikte tools (bij vergelijkende studie, "short-list"),
%% opzetten testopstelling/PoC, uitvoeren testen en verzamelen
%% van resultaten, analyse van resultaten, ...
%%
%% !!!!! LET OP !!!!!
%%
%% Het is uitdrukkelijk NIET de bedoeling dat je het grootste deel van de corpus
%% van je bachelorproef in dit hoofstuk verwerkt! Dit hoofdstuk is eerder een
%% kort overzicht van je plan van aanpak.
%%
%% Maak voor elke fase (behalve het literatuuronderzoek) een NIEUW HOOFDSTUK aan
%% en geef het een gepaste titel.

Het onderzoek wordt uitgevoerd in vier fasen, elk gericht op het beantwoorden van specifieke deelonderzoeksvragen en het realiseren van de doelstellingen van de bachelorproef

\begin{enumerate}
    \item \textbf{Literatuurstudie (week 1-3):} \\
    In een eerste fase wordt een literatuurstudie uitgevoerd om bestaande kennis in kaart te brengen. Hierbij wordt onderzocht welke typen wachtwoordmanagers bestaan en welke info ze bieden, hoe beginnende gebruikers zonder ervaring omgaan met wachtwoordmanagers, welke drempels zij ervaren, welke risico's verbonden zijn aan het niet gebruiken van dergelijke beveiligingstools en welke toegankelijkheidsprincipes relevant zijn voor deze doelgroep. Deze inzichten vormen de basis voor de requirements-analyse.
    \item \textbf{Requirements-analyse (week 3-5):} \\
    In deze fase wordt een requirements-analyse uitgevoerd. Hierbij worden bestaande leeromgevingen rond wachtwoordbeheer geanalyseerd om inzicht te krijgen in hun sterke punten, beperkingen en geschiktheid voor beginnende gebruikers. Op basis hiervan worden de noodzakelijke functionele en niet-functionele requirements voor een veilige webgebaseerde simulatie opgesteld. Deze requirements zullen als referentie dienen voor de proof-of-concept.
    \item \textbf{Technische evaluatie van Unity (week 5-7):} \\
    In een derde fase wordt Unity geëvalueerd op technische haalbaarheid. Hierbij wordt onderzocht in welke mate Unity geschikt is voor webgebaseerde simulatie, specifiek gericht op performance binnen een WebGL-omgeving, toegankelijkheid, UX, beveiliging en ontwikkeltijd. Tevens wordt gekeken naar bestaande frameworks, libraries en design patterns binnen Unity die relevant zijn voor het bouwen van een "fake computer" met fictieve browser en websites. 
    Daarnaast wordt Unity vergeleken met alternatieve webtechnologieën op basis van criteria zoals ontwikkelcomplexiteit, toegankelijkheid, performantie en geschiktheid voor interactieve workshopcontexten.
    \item \textbf{Proof-of-concept ontwikkeling (week 7-12):} \\
    In een vierde fase wordt er een proof-of-concept ontwikkeld. Deze simulatie, gebouwd in Unity met WebGL-export, bevat onder andere volgende functies:
    \begin{itemize}
        \item het aanmaken van fictieve accounts en het genereren en opslaan van wachtwoorden;
        \item feedback- en bewustwordingsmechanismen die gebruikers begeleiden bij veilig wachtwoordbeheer;
        \item een beveiligde sandboxomgeving waarin geen persoonlijke gegevens gebruikt worden.
    \end{itemize} 
    Het proof-of-concept wordt getest door een kleine, representatieve groep beginnende gebruikers uit het netwerk van de student. Het doel is om feedback te verzamelen over de gebruikservaring, bruikbaarheid en eventuele technische beperkingen van de simulatie.
  \end{enumerate}

\textbf{Deliverables per fase:}
\begin{itemize}
    \item Literatuurstudie: overzicht van bestaande kennis, uitdagingen en best practices voor wachtwoordmanagers en educatieve simulaties
    \item Requirements-analyse: document met functionele en niet-functionele requirements voor de webgebaseerde simulatieomgeving
    \item Technische evaluatie: rapport met haalbaarheid, beperkingen en aanbevelingen voor Unity 
    \item Proof-of-concept: werkend webgebaseerd prototype, inclusief observatierapport van gebruikerservaring
\end{itemize}

\lipsum[21-25]

